% 论文三级提纲与写作要点；使用 XeLaTeX 编译。
\documentclass[UTF8,a4paper,zihao=-4]{ctexart}
\usepackage[margin=2.5cm]{geometry}
\usepackage{enumitem}
\setlist{nosep}
\pagestyle{plain}
\clubpenalty=10000\widowpenalty=10000
\ctexset{section={format=\centering\heiti\large,name={,、},number=\chinese{section}},subsection={format=\heiti},subsubsection={format=\heiti}}
\begin{document}
\begin{center}\Large\bfseries 微网与外部电网电力调控策略：论文三级提纲\end{center}
\noindent\textbf{格式依据：}以2024年优秀论文C201为主，依次设置问题重述、问题分析、模型假设、符号说明、数据处理、模型建立与求解、评价改进、参考文献与附录。摘要单页，正文不设目录，控制在30页以内；附录提供完整程序。
\paragraph{摘要（不编号）}约500--800字。以研究目标起笔，依次概述四问的建模思想、算法和真实结果，末段说明信息更新的经济价值及适用条件。关键词建议为线性规划、历史误差场景、滚动优化、储能调度、信息可得性。
\section{问题背景与重述}
\subsection{问题背景}
\subsubsection{微网构成与能量调控目标}说明负载、光伏、储能与外网的关系，界定满足负荷和降低费用两个目标。
\subsubsection{调度尺度与交易机制}说明10分钟时段、日前交易、紧急补电与日内调整的时间顺序。
\subsection{四个问题的核心任务}
\subsubsection{问题1：单日确定性购电与储能调度}强调日初日末储电量相等，以及题目指定表格。
\subsubsection{问题2：全年日前计划与紧急补电}区分全年状态热启动和2--12月输出统计，突出实际偏差结算。
\subsubsection{问题3：光伏预测更新与计划调整}说明四次预报、未执行时段调整和偏差费用。
\subsubsection{问题4：波动电价下的策略比较}区分价格曲线变化与未来价格是否已知，说明两种信息设定。
\section{问题分析}
\subsection{总体建模思路}
\subsubsection{四问之间的继承与递进}以基础物理约束为共同框架，逐步增加运行跨度、预测更新、结算成本与价格信息限制。
\subsubsection{预测、计划、执行与评价的闭环}给出总流程图，明确真实未来数据只参与事后评价。
\subsection{各问题的关键难点}
\subsubsection{问题1：效率、功率与电量边界}说明功率转电量、时间下标、弃电与同时充放电处理。
\subsubsection{问题2：日前未知量与跨日状态}解释缺少日前负荷预测时如何从历史数据构建预测和误差场景。
\subsubsection{问题3：滚动调整与避免重复计费}说明原始计划、调整后计划及紧急购电的区别。
\subsubsection{问题4：电价信息与公平比较}解释可实施策略和价格已知对照，费用比较需同时观察状态差异。
\section{模型假设}
\subsection{物理运行假设}
\subsubsection{时段内功率恒定与损耗口径}明确效率为充电、放电各90\%，其他损耗忽略。
\subsubsection{单向购电与盈余处理}允许弃光或不利用已付费电量，不引入题目未给出的售电收益。
\subsection{决策与结算假设}
\subsubsection{信息可得性与状态观测}历史数据、当前状态与已发布预测可用，未来实际量不可用于可实施计划。
\subsubsection{计划基准与终端边界}仅问题1强制循环边界；其他问题跨日连续。明确两种退购结算解释。
\section{符号说明}
\subsection{索引与输入参数}
\subsubsection{日期、时段与预测发布索引}明确时段1对应0:00--0:10，状态索引为0至144。
\subsubsection{负荷、光伏、电价及设备参数}符号表列出含义与单位，区分预测、实际和场景量。
\subsection{决策变量与评价指标}
\subsubsection{购电、充放电与储电量}统一计划购电、调整购电、紧急补电及弃电符号。
\subsubsection{费用分项与储能利用指标}定义总费用、上调费用、下调费用、等效循环次数。
\section{数据处理与分析}
\subsection{输入数据统一}
\subsubsection{附件字段与时间对齐}核查145个时刻字段对应144个区间，记录模板的10分钟标签错位。
\subsubsection{完整性、非负性与单位检验}报告数据量、缺失值情况、价格范围及能量单位换算。
\subsection{预测与误差场景}
\subsubsection{基于1月份验证集的历史预测选择}比较持久性、移动平均、指数加权与同星期预测，冻结方法后评价2--12月。
\subsubsection{整点预报插值与历史误差构造}避免使用尚未发布的光伏预报，说明场景数量和历史窗口。
\section{模型建立与求解}
\subsection{问题1：单日确定性线性规划模型}
\subsubsection{决策变量、目标函数与物理约束}完整列出费用目标、功率平衡、SOC动态、边界、效率及非负约束。
\subsubsection{线性性质与同时充放电消除}给出允许零成本弃电时的等价变换，说明其成立条件。
\subsubsection{求解流程与指定表格结果}说明稀疏矩阵编码、HiGHS求解、表1与表2中的真实数据。
\subsubsection{费用与储能策略分析}结合日内曲线解释低价购电、高价放电及效率造成的损耗。
\subsection{问题2：含紧急补电风险的日前调度模型}
\subsubsection{历史误差场景与两阶段费用代理}设置共同日前计划与场景缺额变量，将5倍价格的惩罚纳入目标。
\subsubsection{实际运行反馈与跨日连续性}给出安全裁剪、补电、弃电和状态更新算法，说明计划状态与实际状态的差别。
\subsubsection{典型日与统计期计算结果}给出四个指定日期的表1--3及全年费用分解，说明result2.xlsx结构。
\subsubsection{紧急购电成因与基准比较}比较无风险代理的确定性日前模型，分析计划多购与紧急补电之间的权衡。
\subsection{问题3：考虑调整费用的滚动优化模型}
\subsubsection{原始计划、调整量与偏差变量}以0:00计划为统一结算基准，明确每个时段只结算最终生效计划。
\subsubsection{滚动目标函数与执行算法}给出主结算口径与退费口径的目标函数，按0、6、12、18时启动优化。
\subsubsection{典型日结果与费用分解}展示原计划、调整后购电、紧急购电及储能实际轨迹。
\subsubsection{更新频率与结算口径对比}比较每日一次、中午更新、四次更新，解释费用改善不必然相对问题2成立。
\subsection{问题4：波动电价下的调控模型}
\subsubsection{未来电价未知的可实施模型}以历史价格预测作为优化输入，以实际价格计算全部交易费用。
\subsubsection{日前价格已知的对照模型}保持负荷和光伏信息不变，单独改变价格信息可得性。
\subsubsection{计算结果与储能经济性比较}给出两类结果表、典型日调度、紧急电量和成本构成。
\subsubsection{稀疏模型规模与运行时间}报告变量、约束和求解次数，不把经验耗时误写为普遍复杂度界。
\subsection{模型检验与讨论}
\subsubsection{物理约束及费用核验}检查逐时平衡、能量动态、容量、功率、跨日衔接与费用相加。
\subsubsection{场景窗口与结算解释敏感性}报告7、14、28天窗口及退费口径下的实际费用变化。
\section{模型评价与改进}
\subsection{模型评价}
\subsubsection{模型优点}总结统一物理模型、计算效率、信息时序清晰和结果可核验。
\subsubsection{模型局限}说明固定储能计划的风险代理与实际反馈并不完全相同、缺少电池衰减和全局最优保证。
\subsection{改进方向}
\subsubsection{更准确的多阶段不确定性建模}引入场景相关性、储能分场景补救及分布鲁棒控制。
\subsubsection{长期寿命与预测改进}结合循环寿命成本、长期终端价值和外生气象信息。
\section{参考文献}
\noindent 引用真实可核查的线性规划、滚动控制、预测评价和求解器文献。范文仅作为格式参考，不移用其研究结论。
\section*{附录说明}
\subsection*{附录A：支撑材料清单与复现步骤}列出源码、输入要求、5个结果表和编译命令。
\subsection*{附录B：全部可运行程序}完整收录预测、模型构建、全年仿真、结果导出与图表生成程序。
\subsection*{附录C：典型日完整紧急购电表}篇幅较长的逐时事件表在此补充，正文保留关键表格与解释。
\end{document}
